\documentclass[10pt]{article}

\usepackage[utf8]{inputenc}
\usepackage[T1]{fontenc}
\usepackage[spanish]{babel}

\usepackage{aleph-comandos}
\usepackage{moodle}
% Se recomienda leer la documentación de estos
% paquetes en https://www.alephsub0.org/recursos/

% -- Paquetes adicionales
\usepackage{amssymb}
\usepackage{mathpazo}

% -- Datos
\title{Reactivos Álgebra Lineal}
\author{Evelyn Cueva}

\usepackage[margin=2cm]{geometry}

\begin{document}
\maketitle

\begin{quiz}{Algebra-Lineal-Examen-Media-Carrera-2024B}

% Ejercicio de cálculo para evaluar núcleo e imagen de una transformación lineal. Matriz transpuesta y composición de transformaciones 

\begin{multi}{algebra-lineal-2024B-01}
    Considere $T$ una tranformaci\'on lineal que transpone cualquier matriz de $2\times 2.$
    Indique cu\'al de las siguientes afirmaciones es falsa:
    \item[fraction=0] $T\circ T$ es la transformaci\'on identidad. 
    \item[fraction=0] $T(A)=\begin{pmatrix}
        0&0\\
        0&0
    \end{pmatrix}$ si y solo si $A$ es la matriz nula.
    \item[fraction=0] La dimensión del espacio imagen de $T$ es 4. 
    \item* Si $M$ es una matriz de $2\times 2$, $T(M)=-M$ no puede ocurrir. 
\end{multi}



% Ejercicio de cálculo para evaluar propiedades de determinantes.

\begin{multi}{algebra-lineal-2024B-02}
Si $A$ y $B$ son dos matrices de $3\times 3$ tales que $|A| = 4$ y $|B|=3$, donde $|A|$ denota el determinante de $A$. ¿Cu\'al es el determinante de $A^{-1}B$?
     \item* $\frac 34$
     \item[fraction=0]  12
     \item[fraction=0] 7
     \item[fraction=0] $\frac 43$
\end{multi}

% Pregunta conceptual y de calculo para evaluar independencia lineal y espacios generados por listas de diferentes tamaños.

\begin{multi}{algebra-lineal-2024B-03}
De los siguientes conjuntos de vectores en $\Rbb^3$, determine cu\'al no genera todo el espacio $\R^3$.
    \item[fraction=0]  $\left\lbrace(1,-1,1),(-1,-1,1),(0,1,1)\right\rbrace$
    \item[fraction=0]  $\left\lbrace (1,-1,1),(-1,-1,1),(0,1,1),(0,-1,3) \right\rbrace$
    \item*  $\left\lbrace(8,-3,4),(0,3,0),(8,0,4)\right\rbrace$
    \item[fraction=0] $\left\lbrace (1,0,0),(0,1,0),(0,0,1) \right\rbrace$
\end{multi}

% Pregunta conceptual para identificar la relacion entre propiedades de matriz inversa, sistemas de ecuaciones y sistemas homogeneos

\begin{multi}{algebra-lineal-2024B-04}
Sea $A$ una matriz de tama\~no $n\times n$. De las siguientes afirmaciones, seleccione aquella que no sea equivalente con las dem\'as:
    \item[fraction=0] $A$ es no singular.
    \item[fraction=0] El sistema $Ax=0$ tiene soluci\'on \'unica trivial.
    \item* $A^{-1}$ no existe. 
    \item[fraction=0] El determinante de $A$ es no nulo.
\end{multi} 
\end{quiz}

\end{document}